% =========================================================
% Order-Preserving Maps, Embeddings, and Isomorphisms
% =========================================================

\section{Maps Between Ordered Sets}
\label{sec:maps-ordered-sets}

\begin{remark*}[Toolkit: Maps Between Ordered Sets]
The structure-preserving maps between posets form a hierarchy of increasing
strength.  An order-preserving map carries the direction of the order
forward.  An order-embedding carries it in both directions.  An
order-isomorphism is a surjective order-embedding --- a perfect structural
identification between two posets.
\end{remark*}

% ---------------------------------------------------------
% def:order-preserving-map
% ---------------------------------------------------------

\begin{definitionbox}{Definition (Order-Preserving Map)}
\begin{definition}[Order-Preserving Map]
\label{def:order-preserving-map}
Let $(P, \leq_P)$ and $(Q, \leq_Q)$ be partially ordered sets.  A map
$\varphi : P \to Q$ is \emph{order-preserving} (or \emph{monotone}) if
\[
\forall x, y \in P, \quad
x \leq_P y \;\Longrightarrow\; \varphi(x) \leq_Q \varphi(y).
\]
\end{definition}
\end{definitionbox}

\begin{remark*}[Standard quantified statement]
\[
\operatorname{OrderPreserving}(\varphi, P, Q)
\;\coloneqq\;
\forall x, y \in P,\;
x \leq_P y \Rightarrow \varphi(x) \leq_Q \varphi(y).
\]
\end{remark*}

\begin{remark*}[Negated quantified statement]
\[
\neg\operatorname{OrderPreserving}(\varphi, P, Q)
\;\coloneqq\;
\exists x, y \in P,\;
x \leq_P y \land \varphi(x) \not\leq_Q \varphi(y).
\]
\end{remark*}

\begin{remark*}[Interpretation]
An order-preserving map carries the direction of the order from $P$ into
$Q$: whenever $x$ is at or below $y$ in $P$, $\varphi(x)$ is at or below
$\varphi(y)$ in $Q$.  The implication is one-directional --- order-preserving
maps need not reflect the order back.  Two order-preserving maps compose
to an order-preserving map, making order-preserving maps the morphisms
of the category of posets.
\end{remark*}

\begin{dependencies}
\hyperref[def:lattice-order-partial-order]{Partial order}.
\end{dependencies}

% ---------------------------------------------------------
% def:lattice-order-embedding
% ---------------------------------------------------------

\begin{definitionbox}{Definition (Order-Embedding)}
\begin{definition}[Order-Embedding]
\label{def:lattice-order-embedding}
Let $(P, \leq_P)$ and $(Q, \leq_Q)$ be partially ordered sets.  A map
$\varphi : P \to Q$ is an \emph{order-embedding} if
\[
\forall x, y \in P, \quad
x \leq_P y \;\Longleftrightarrow\; \varphi(x) \leq_Q \varphi(y).
\]
When $\varphi$ is an order-embedding, one writes $\varphi : P \hookrightarrow Q$.
\end{definition}
\end{definitionbox}

\begin{remark*}[Standard quantified statement]
\[
\operatorname{OrderEmbedding}(\varphi, P, Q)
\;\coloneqq\;
\forall x, y \in P,\;
x \leq_P y \Longleftrightarrow \varphi(x) \leq_Q \varphi(y).
\]
\end{remark*}

\begin{remark*}[Interpretation]
An order-embedding both preserves and reflects the order: the order
structure of $P$ is faithfully copied into $Q$ without distortion.
Unlike mere preservation, reflection means that order-relations between
images in $Q$ genuinely originate from order-relations in $P$.  Every
order-embedding is injective
(\hyperref[lem:order-embedding-injective]{Lemma~\ref*{lem:order-embedding-injective}})
and identifies $P$ order-isomorphically with its image in $Q$.
\end{remark*}

\begin{dependencies}
\hyperref[def:order-preserving-map]{Order-preserving map}.
\end{dependencies}

% ---------------------------------------------------------
% def:order-isomorphism
% ---------------------------------------------------------

\begin{definitionbox}{Definition (Order-Isomorphism)}
\begin{definition}[Order-Isomorphism]
\label{def:order-isomorphism}
An \hyperref[def:lattice-order-embedding]{order-embedding} $\varphi : P \to Q$
that is surjective is an \emph{order-isomorphism}.  When an
order-isomorphism exists, $P$ and $Q$ are said to be
\emph{order-isomorphic}, written $P \cong Q$.
\end{definition}
\end{definitionbox}

\begin{remark*}[Standard quantified statement]
\[
\operatorname{OrderIso}(\varphi, P, Q)
\;\coloneqq\;
\operatorname{OrderEmbedding}(\varphi, P, Q)
\;\land\;
\forall q \in Q,\; \exists p \in P,\; \varphi(p) = q.
\]
\end{remark*}

\begin{remark*}[Interpretation]
An order-isomorphism is a bijection that perfectly translates the order
of $P$ into $Q$ and back.  Two order-isomorphic posets are
order-theoretically identical: any statement about $P$ expressed purely
in terms of $\leq_P$ translates to an equally valid statement about $Q$
under $\varphi$.  Order-isomorphism is the correct notion of equality
for posets.
\end{remark*}

\begin{dependencies}
\hyperref[def:lattice-order-embedding]{Order-embedding}.
\end{dependencies}

% ---------------------------------------------------------
% lem:order-embedding-injective
% ---------------------------------------------------------

\begin{lemmabox}{Lemma}
\begin{lemma}[Every Order-Embedding Is Injective]
\label{lem:order-embedding-injective}
Let $(P, \leq_P)$ and $(Q, \leq_Q)$ be partially ordered sets, and let
$\varphi : P \to Q$ be an order-embedding.  Then $\varphi$ is injective.
\hyperref[prf:order-embedding-injective]{Go to proof.}
\end{lemma}
\end{lemmabox}

\begin{remark*}[Standard quantified statement]
\[
\operatorname{OrderEmbedding}(\varphi, P, Q)
\;\Longrightarrow\;
\forall x, y \in P,\; \varphi(x) = \varphi(y) \Rightarrow x = y.
\]
\end{remark*}

\begin{remark*}[Interpretation]
Injectivity follows automatically from the biconditional in the embedding
definition.  If $\varphi(x) = \varphi(y)$, the image is simultaneously
$\leq$ and $\geq$ itself in $Q$; order-reflection pulls this back to $P$,
and anti-symmetry then forces $x = y$.
\end{remark*}

\begin{dependencies}
\hyperref[def:lattice-order-embedding]{Order-embedding},
\hyperref[def:lattice-order-partial-order]{anti-symmetry of partial orders}.
\end{dependencies}
